\subsection*{Pooling probability on identical physical spans}
For each pictured question and passage, the generic reader and three adapted endpoint heads propose four strips, $v_1,\ldots,v_4$. The selector combines the shared contextual summary with each head's top joint score, first-to-second score gap, posterior mass, start position and span length. It also receives six pairwise bitmap agreements, the predicted relation confidence and the pictured question length. It has 42,340 trainable parameters. Its output weights $\pi_i$ sum to one. For a proposed image $y$, let $P(y)=\sum_{i:v_i=y}\pi_i$. The prediction maximizes $P(y)$ over proposed strips, breaking ties by original proposal order. Training minimizes the negative logarithm of the total probability of all accepted images whenever any is available. Of 15,704 training questions, 12,565 supplied such a gradient. With nonnegative integer weights of total $W$, Lean verifies that an incorrect pooled maximizer implies $W\leq 2(W-W_y)$ for any accepted reference of weight $W_y$. A common denominator extends this sufficient condition to rational weights. The proof supplies neither a convergence result nor a guarantee of generalization.

The fit mixed 7,000 original CMRC and DRCD training questions and 8,704 pictured exchanges. Evaluation used the first 416 passage cells and required an annotated span of at most 48 cells. The best epoch maximized the mean of natural and exchange validation accuracies on 512 natural questions and 896 exchanges. The validation sources were withheld from selector training; the frozen participant heads had seen the earlier exchange training pool. Two fits used matched initialization and 15 epochs, with label vectors shuffled across questions in the control. Processing training and validation images took 81.0 seconds and fitting both heads took 4.8 seconds on one RTX 4090 D. Neither time includes backbone training. The new 384-question evaluation used one question per previously excluded paragraph from CMRC development or DRCD test. Upstream exposure is unverified. Accuracy requires exact packed glyphs. Complete factual tests also require termination in both requested word orders. For the separate availability measurement, ties among spans follow physical index order. We compared all proposed images to all accepted references without changing weights or generated answers.

